# Finite Reservoir Quorum (FRQ): fitted coordinates and likelihood

This note is the short live specification of the FRQ model. Derivations,
the classic-counter submodels, error signatures and identification results
are in `CounterAndQuorumModels.tex`. The sampler-facing coordinates are

\[
  (K,\;v,\;sv,\;h,\;t_0,\;\pi_{lo},\;\mathrm{clock}),
\]

with \(\pi_{lo}\) (`plo`) and `clock` as optional extensions. Every coordinate is a
mechanism. None is a response-time summary.

## 1. Construction

A reservoir of \(N\) units has each unit available with probability \(p\).
An available unit registers after a delay with CDF \(G\), and the accumulator
responds when \(K\) units have registered. With \(\alpha=K\),
\(\beta=N-K+1\) (both real and positive), and \(q(u)=pG(u)\),

\[
  F(u)=I_{q(u)}(\alpha,\beta),\qquad
  f(u)=p\,g(u)\,\frac{q^{\alpha-1}(1-q)^{\beta-1}}{B(\alpha,\beta)},\qquad
  S(u)=I_{1-q(u)}(\beta,\alpha).
\]

The registration law has rate \(\lambda\) and shape \(c\) (parameter `clock`):

\[
  G(u)=1-\left(1+c^2\lambda u\right)^{-1/c^2},
\]

which becomes \(G(u)=1-e^{-\lambda u}\) at \(c=0\). In every case
\(g(0)=\lambda\). The observed time is \(T=t_0+U\). The eventual completion
probability is \(h=F(\infty)=I_p(\alpha,\beta)\), and \(U=\infty\) with
probability \(1-h\).

## 2. Counter reading

With \(\phi(u)=q/(1-q)\),

\[
  F(u)=P\{\mathrm{Gamma}(\alpha,1)\le W\phi(u)\},\qquad
  W\sim\mathrm{Gamma}(\beta,1).
\]

Each accumulator is therefore a Poisson counter with threshold \(K\), a trial
intensity \(W\) with CV \(1/\sqrt\beta\), and an operational clock \(\phi\).
The clock satisfies \(\phi'(0)=p\lambda\) and saturates at \(p/(1-p)\) when
\(p<1\), which is the evidence ceiling. At \(p=1\) the clock accelerates for
\(c<1\), is exactly linear at \(c=1\), and decelerates for \(c>1\).

## 3. Fitted chart

\[
  K=\alpha,\qquad v=\beta p\lambda,\qquad sv=v\,\beta^{-1/2},\qquad h=F(\infty).
\]

- \(K\) is the event threshold.
- \(v\) is the mean initial input rate in events per second.
- \(sv\) is the between-trial SD of the input rate. Its CV, \(cv=sv/v\), is
  derived. \(sv>v\) is allowed.
- \(h\) is the probability that the accumulator ever responds.

The kernel inverts the chart as

\[
  \beta=(v/sv)^{2},\qquad p=Q_{\alpha,\beta}(h),
  \qquad \lambda=\frac{v}{\beta p}.
\]

With \(\pi_{lo}>0\), \(p\) is instead the root of the mixed completion
probability (Section 4). Either way it depends only on \((K,sv/v,h,\pi_{lo})\)
and is cached on those
four arguments. A per-trial \(v\) therefore costs a quantile per distinct value
only when \(sv>0\) and \(h<1\). At \(h=1\), \(p=1\) exactly. Below one, \(1-p\)
is solved from its own quantile, \(Q_{\beta,K}(1-h)\), so the
defect stays exact when \(\beta<1\) puts \(p\) within rounding of one. The
tails at \(u=\infty\) return \(h\) and \(1-h\) directly. \(1-p\) and
\(1-q(u)\) are held as logarithms, since for \(\beta<1\) they underflow while
\((1-q)^\beta\) is still of order one.

**The cv chart.** `FRQ(parameterization = "cv")` samples
\(cv=sv/v=\beta^{-1/2}\) in `sv`'s place, with `sv`'s link, bound, default
and exception. The charts are the same model (equal likelihood at
\(sv=cv\,v\)); \(cv=0\) is the \(sv=0\) branch. With
\((K,cv,h,\pi_{lo},\mathrm{clock})\) fixed, \(v\) is a pure time scale,
\(F(u;v)=F(vu;1)\), and the quantile key is free of \(v\). In the `sv`
chart this holds only along \(sv\propto v\). The mapping layer then reports
`sv` in place of `cv`.

**The \(sv=0\) branch** is exact, not a limit approximation. The kernel also
takes it whenever \(cv<10^{-4}\), where the beta branch is within about
\(5\times10^{-9}\) of the limit. Let \(P_K\) be
the regularised lower incomplete gamma function.

- At \(h=1\): \(F(u)=P_K(vu)\), the homogeneous Poisson counter, and the
  clock drops out.
- At \(h<1\): let \(E=P_K^{-1}(h)\) be the expected total
  evidence. Then the rate is \(\lambda=v/E\) and \(F(u)=P_K(E\,G(u))\).

**Clock and rate variability.** `clock` and `sv` carry the same contagion
signature. At \(K=1\) only \(\beta/c^2\) enters the law, so they are exactly
confounded. At \(K=3\) they are distinct in principle, but with both free their
measured posterior correlation is 1. The constructor therefore fixes `clock = 1`, the constant-rate counter in which
\(sv\) is exactly gamma-mixed rate variability. `clock = 0` is the
accelerating variant, which is exactly the linear self-exciting counter with
\(\gamma=\lambda\) and \(\nu=v\).

**Rate SD per accumulator.** Designs use `sv ~ lM`. Because `K` counts
events, `v` and `sv` are absolute and no `sv` needs fixing; at h = 1 both are
identified (per-trial SD of log sv about 6-15 for the correct accumulator), but
correlate 0.9-0.98 with their own `v`. The Gamma rate makes the sign of the
error effect depend on `sv_e / sv_c`: equal CVs give slow errors, equal SDs
give fast errors at large variability (see `CounterAndQuorumModels.tex`).

## 4. Low-threshold state \(\pi_{lo}\)

On each trial each accumulator independently commits at the low threshold
\(K_{lo}=1\) with probability \(\pi_{lo}\) (`plo`), and at \(K\) otherwise.
Both thresholds share \(\beta\), \(p\) and the clock, so

\[
  F(u)=(1-\pi_{lo})\,I_{q(u)}(K,\beta)+\pi_{lo}\,I_{q(u)}(1,\beta),
  \qquad I_q(1,\beta)=1-(1-q)^\beta,
\]

and likewise for the density, the survivor and the \(sv=0\) branch with
\(P_1(X)=1-e^{-X}\). In the counter reading the \(\mathrm{Gamma}(K,1)\) level is
\(E_1+E_r\), and the low state finishes at \(E_1\), the first event.
\(K_{lo}\) is a kernel constant (`FRQ_K_LO`), not a column.

- The completion probability is
  \(h=(1-\pi_{lo})I_p(K,\beta)+\pi_{lo}I_p(1,\beta)\). At \(h=1\), \(p=1\).
  Below one, \(p\) is the root, bracketed by the two pure quantiles and solved
  by safeguarded Newton in \(\operatorname{logit}p\) (in \(\log E\) at
  \(sv=0\)).
- For \(K>1\), \(f(0^+)=\pi_{lo}v\): the leading edge is first order in time,
  and the density jumps at \(t_0\).
- With only \(v\) differing between accumulators, the low state gives fast
  errors; every \(\pi_{lo}=0\) member gives slow ones.
- At \(K=1\) the components coincide and \(\pi_{lo}\) drops out.
- \(\pi_{lo}=0\) is the base model and the default. The draw is per
  accumulator, so the race still factors.

## 5. Contract

| Parameter | Link    | Bound                  | Default | Exception |
|-----------|---------|------------------------|---------|-----------|
| `K`       | 1 + exp | [1, Inf)               | 3       | 1         |
| `v`       | log     | [1e-4, Inf)            | 6       |           |
| `sv`      | log     | [1e-4, Inf)            | 0       | 0         |
| `cv`      | log     | [1e-4, Inf)            | 0       | 0         |
| `h`       | probit  | [1e-6, 1]              | 1       | 1         |
| `t0`      | log     | [0.05, Inf)            | 0       | 0         |
| `plo`     | probit  | [1e-6, 1]              | 0       | 0         |
| `clock`   | log     | [1e-4, Inf)            | 1       | 0         |

`cv` replaces `sv` under `parameterization = "cv"`.

The bounds are numerical guards, with one exception: the `t0` floor is a
psychological floor.

- \(K\ge1\) is required because for \(K<1\) the density is singular at
  \(t_0\).

The mapping layer reports `beta`, `p`, `lambda`, `Etot` and `cv`
(\(sv/v\)). `Etot` is \(\beta p/(1-p)\), or \(E\) at \(sv=0\), and is infinite
at \(h=1\).

Not covered by the kernel:

- rate variability and self-excitation together;
- the geometric threshold excess \(\omega\), which rate CV mimics while being
  better identified;
- rates correlated across accumulators;
- the log-odds threshold frailty \(\delta\), which \(\pi_{lo}\) replaced:
  it is confounded with \(K\) and gives slow errors (see
  `CounterAndQuorumModels.tex`).

The removed `PCOUNTERcorr` model is described in
`PCOUNTER_correlated_rates.md` and can be recovered from git at `75582521`.
